\section{Estratégias de crescimento}
\subsection{Estratégia 1 -- Densidade local por parcerias (principal)}
Concentrar a operação inicial em uma área de Belo Horizonte ou região metropolitana que possa ser acompanhada presencialmente. A aquisição de tutores será articulada por clínicas, ONGs e lideranças comunitárias. O lado tutor permanece gratuito no núcleo; a monetização é testada com clínicas, assinaturas opcionais e patrocínios. A escolha respeita o problema de arranque bilateral: a equipe precisa garantir oferta local suficiente para que a aquisição de tutores gere valor concreto.

\textbf{Componentes:}
\begin{enumerate}[label=\alph*)]
\item selecionar 2--3 clínicas e até 2 organizações comunitárias para um piloto, sem afirmar parceria antes de contrato/aceite formal;
\item usar QR individualizado por parceiro, redes sociais e indicação, com registro da origem de cada cadastro;
\item medir a conversão de cadastro em ativação e encaminhamento concluído;
\item abrir nova praça somente após os gatilhos da Seção 7.
\end{enumerate}

\subsection{Estratégia 2 -- Aumentar contribuição por usuário}
Testar a disposição a pagar pelo PetCuida+ e pelo Plano Parceiro, sem lançar serviços adicionais antes de comprovar demanda e capacidade de entrega. A equipe deve comparar intenção declarada com comportamento real: teste de preço, pré-adesão ou contrato piloto, seguido de observação de conversão e cancelamento. Patrocínio de campanhas pode reduzir o custo do fundo de créditos, mas só entra na projeção quando houver compromisso formal e regras de repasse definidas.

\subsection{Piloto territorial assistido de oito semanas}
\begin{table}[H]\centering
\caption{Desenho do piloto e decisões}
\begin{tabularx}{\textwidth}{L{2cm}L{2.5cm}Y}
\toprule
Período & Foco & Atividades e evidências\\\midrule
Semanas 1--2 & Descoberta & 15--20 entrevistas com tutores; conversas com 5 clínicas e 2 organizações; mapear barreiras, frequência de uso e oferta disponível.\\
Semanas 3--4 & Usabilidade & 10 testes de tarefas sem ajuda; registrar taxa de conclusão, erros e dúvidas; revisar QR, créditos e linguagem da triagem.\\
Semanas 5--7 & Piloto concierge & Até 40--50 tutores e 2--3 clínicas, se houver consentimento, capacidade e salvaguardas; acompanhar cadastro, lembretes e encaminhamentos manualmente.\\
Semana 8 & Decisão & Continuar, iterar ou interromper; apresentar contagens absolutas, taxas e limitações, sem extrapolar a amostra.\\
\bottomrule
\end{tabularx}\fonteautores
\end{table}

O piloto não deve realizar diagnóstico veterinário automatizado. A triagem será informativa e orientativa, com regras revisadas por profissional habilitado e encaminhamento claro para atendimento presencial. Antes da coleta de dados reais, a equipe deve concluir avaliação de privacidade, controle de acesso, minimização, retenção, backup, resposta a incidentes e contratos com operadores, conforme a LGPD e orientações da ANPD \cite{brasil2018,anpd2021,anpd2022}.